% from ega4-2-fr.tex
\oldpage[IV-2]{148}

\item[(ii)] Comme on l'a rappelé en (i), il résulte de (4.6.1) que,
lorsque $X$ est un préschéma algébrique sur $k$ qui n'est pas
géométriquement réduit, il y a des extensions radicielles finies $k'$ de
$k$ telles que $X'=X\otimes_k k'$ soit non réduit. Il est intéressant de
donner un exemple de ce fait où $X$ est un schéma régulier sur $k$, tel
que $k$ soit algébriquement fermé dans le corps $K$ des fonctions
rationnelles de $X$. Soit $k$ un corps de caractéristique $p>0$, dans
lequel il existe deux éléments $a$, $b$ formant une famille $p$-libre sur
$k^p$. Désignant encore par $B$ l'anneau $k[S,T]$, considérons cette fois
le polynôme $P(S,T)=T^p-aS^p-b$; comme $aS^p+b$ n'est pas une puissance
$p$-ième dans $k(S)$, $P$ est irréductible comme polynôme de $k(S)[T]$,
et le schéma $X_0=\operatorname{Spec}(B/PB)$ est donc une courbe affine
intègre sur $k$, dont le corps des fonctions rationnelles est
$K=k(S)[t]$, où $t$ est une racine de $P$; montrons que $k$ est
algébriquement fermé dans $K$. Supposons en effet que $K$ contienne un
élément $z$ algébrique sur $k$ et non dans $k$; on aurait alors aussi
$z\notin k(S)$, donc $[k(S)[z]:k(S)]>1$; comme $[K:k(S)]=p$, on aurait
$z^p\in k(S)$, donc $c=z^p\in k$ puisque $z$ est algébrique sur $k$;
mais on aurait $t^p=aS^p+b\in k^p(S^p)(c)$, donc $a$ et $b$
appartiendraient à $k^p(c)$, ce qui est absurde. Soit alors $X$ la
normalisée de la courbe $X_0$ dans $K$, qui est donc une courbe normale
(et par suite régulière) sur $k$. Si $k'=k(a^{1/p},b^{1/p})$, il est
clair que $aS^p+b$ est une puissance $p$-ième dans $k'(S)$, donc
$K\otimes_k k'$ n'est pas réduit, ni a fortiori le schéma
$X\otimes_k k'$.
\end{enumerate}
\end{remarks}

\begin{definition}[6.7.6]
\label{IV.2.s6.7.6-fr}
Soient $k$ un corps, $X$ un $k$-préschéma localement noethérien. On dit
que $X$ est géométriquement régulier en un point $x$ (resp. possède la
propriété $(R_n)$ géométrique en un point $x$, resp. est géométriquement
normal en un point $x$) si, pour toute extension finie $k'$ de $k$,
$X'=X_{(k')}=X\otimes_k k'$ est régulier (resp. possède la propriété
$(R_n)$, resp. est normal) en tout point $x'$ dont la projection dans
$X$ est $x$. On dit que $X$ est géométriquement régulier (resp. possède
la propriété $(R_n)$ géométrique (ou encore est géométriquement régulier
en codimension $\leqslant n$), resp. est géométriquement normal) si $X$
est géométriquement régulier (resp. possède la propriété $(R_n)$
géométrique, resp. est géométriquement normal) en tout point.

On dit qu'une algèbre $A$ sur $k$ est un anneau géométriquement régulier
(resp. géométriquement normal, resp. géométriquement réduit, resp.
possédant la propriété $(R_n)$ géométrique) si $\operatorname{Spec}(A)$
a la même propriété.

Si $X=\operatorname{Spec}(K)$, où $K$ est une extension de $k$, il revient
au même de dire que $X$ est géométriquement régulier, ou géométriquement
normal, ou géométriquement réduit, ou que $K$ est une extension séparable
de $k$: cela résulte de (4.6.1) et du fait que si $K$ est une extension
séparable de $k$ et $k'$ une extension finie de $k$, $K\otimes_k k'$ est
composé direct d'un nombre fini de corps (Bourbaki, Alg., chap. VIII,
§~7, n°~3, cor. 1 du th. 1).
\end{definition}

